% !TEX encoding = UTF-8 Unicode

% ------------------------------------------------------------------------------------------------------
%	Formatvorlage für wissenschaftliche Arbeiten (Diplomarbeit, Bachelorarbeit, Masterarbeit, Seminararbeit, Projektarbeit, Dokumentation)
% ------------------------------------------------------------------------------------------------------
%	ursprünglich erstellt von Stefan Macke, 24.04.2009
%	http://blog.stefan-macke.de
%
%	erweitert von Felix Rupp
%	http://www.felixrupp.com/
%
%
%
% erweitert und angepasst von Moritz Rupp
% moritz.rupp@gmail.com
%
% Version 1.0
% Datum: 01.06.2013


% Dokumentenkopf ---------------------------------------------------------------------------------------
%   Diese Vorlage basiert auf "scrreprt" aus dem koma-script.
% ------------------------------------------------------------------------------------------------------
\documentclass[
    12pt, % Schriftgröße
    DIV10, % Änderung der Größe des Satzspiegels (bedruckbarer Bereich einer Seite), nur in Verbindung mit koma-script verwendbar
    UKenglish, % für Umlaute, Silbentrennung etc.
    a4paper, % Papierformat
    oneside, % einseitiges Dokument
    titlepage, % es wird eine Titelseite verwendet
    parskip=half, % Abstand zwischen Absätzen (halbe Zeile)
    headings=normal, % Größe der Überschriften verkleinern
    listof=totoc, % Verzeichnisse im Inhaltsverzeichnis aufführen
    bibliography=totoc, % Literaturverzeichnis im Inhaltsverzeichnis aufführen
    index=totoc, % Index im Inhaltsverzeichnis aufführen
    captions=tableheading, % Beschriftung von Tabellen unterhalb ausgeben
    final % Status des Dokuments (final/draft)
]{scrreprt}

% UTF8 und T1 Fontencoding -----------------------------------------------------------------------------
\usepackage[utf8]{inputenc}
\usepackage[T1]{fontenc}


% Meta-Informationen -----------------------------------------------------------------------------------
%   Informationen über das Dokument, wie z.B. Titel, Autor, Matrikelnr. etc
%   werden in der Datei Meta.tex definiert und können danach global
%   verwendet werden.
% ------------------------------------------------------------------------------------------------------
% !TEX encoding = UTF-8 Unicode
% !TEX root =  Paper.tex

% Meta information ------------------------------------------------------------------------------------
% Definiton of global parameters, they can be used in the whole document (e.g. on the titlepage)
% -------------------------------------------------------------------------------------------------------
\newcommand{\myTitle}{The Effects of National Business Systems on HRM in Multinational Companies} % the title for text
\newcommand{\myTitleTitlepage}{The Effects of National Business Systems on HRM in \\ Multinational Companies} % the title for the titlepage

\newcommand{\mySubtitle}{} 
\newcommand{\mySubtitleTitlepage}{}

\newcommand{\study}{General Management}
\newcommand{\area}{Faculty of Economic and Social Sciences}

\newcommand{\myAuthor}{Michaela Wunder}

\newcommand{\keywords}{}

\newcommand{\born}{Ellwangen}

\newcommand{\firstreviewer}{Professor Dr.\,Kerstin Pull}
\newcommand{\secondreviewer}{}

\newcommand{\submissionDate}{09/04/2015}

\newcommand{\university}{Eberhard Karls Universität}
\newcommand{\place}{Tübingen}

\newcommand{\creator}{texmaker 4.2}

\newcommand{\bibliographyheading}{Reference list}

\newcommand{\TableCellPadding}{1.2}


% benötigte Packages -----------------------------------------------------------------------------------
%   LaTeX-Packages, die benötigt werden, sind in die Datei Packages.tex
%   "ausgelagert", um diese Vorlage möglichst übersichtlich zu halten.
% ------------------------------------------------------------------------------------------------------
% !TEX encoding = UTF-8 Unicode
% !TEX root =  Paper.tex

% Customisation of the page layout ---------------------------------------------------
%   see Layout.tex
% -----------------------------------------------------------------------------------------
\usepackage[
    automark, % automatically create chapter name in the headline
    headsepline, % seperation line beneath headline
    ilines % left seperate line
]{scrpage2}

% Hacks in order to prevent warnings
\usepackage{scrhack}

% Customisation of the language -------------------------------------------------
\usepackage[british]{babel}


% Special character 
% -----------------------------------------------------------------------------------------
\usepackage{textcomp} % Euro-sign etc.


% Font --------------------------------------------------------------------------------
\usepackage{lmodern} % better fonts
\usepackage{relsize} % font size relatively


% Better underlining ---------------------------------------------
\usepackage[normalem]{ulem}


% Graphics -----------------------------------------------------------------------------
% Enable the inclusion of JPEG
\usepackage[dvips,final]{graphicx}
% Path to the graphics
\graphicspath{{Images/}}
\usepackage{picture}


% Commands from AMSTeX for mathematical symbols, e.g.  \boldsymbol \mathbb
\usepackage{amsmath,amsfonts}


% Use Euro sign ------------------------------------------------------------
\usepackage{eurosym}


% For index with \printindex -----------------------------------------------
\usepackage{makeidx}


% Easy definition of margins and line spacing etc. ------------
\usepackage{setspace}
\usepackage{geometry}

% Floating figures ---------------------------------------------------------
\usepackage[vflt]{floatflt}
\usepackage{subfig}
\usepackage{tocbasic}

% Inclusion of program code -----------------------------------------------
\usepackage{listings}
\usepackage{xcolor} 
\definecolor{hellgelb}{rgb}{1,1,0.9}
\definecolor{colKeys}{rgb}{0.8,0,0.5}
\definecolor{colIdentifier}{rgb}{0.6,0,0.3}
\definecolor{colComments}{rgb}{0,0.5,0}
\definecolor{colString}{rgb}{0,0,1}

\lstset{
    float=htbp,
    basicstyle=\ttfamily\color{black}\small\smaller,
    identifierstyle=,%\color{colIdentifier},
    keywordstyle=\color{colKeys}\bfseries,
    stringstyle=\color{colString},
    commentstyle=\color{colComments},
    columns=flexible,
    tabsize=4,
    frame=single,
    extendedchars=true,
    showspaces=false,
    showstringspaces=false,
    numbers=left,
    numberstyle=\tiny,
    breaklines=true,
    backgroundcolor=\color{hellgelb},
    breakautoindent=true,
    escapeinside={(*}{*)},
    literate={Ö}{{\"O}}1 {Ä}{{\"A}}1 {Ü}{{\"U}}1 {ß}{{\ss}}2 {ü}{{\"u}}1 {ä}{{\"a}}1 {ö}{{\"o}}1 {µ}{\textmu}1
 }

% --------------------------------------------------------------------------------------------
% 
% Definition of source code
%
% Sample Defininition of CSS:
\lstdefinelanguage{CSS}{
    morestring=[b]',
    morestring=[b]",
    comment=[l]{/*}{*/},
    sensitive=false,
    morekeywords={accelerator,adjust,after,align,attachment,azimuth,background,before,behavior,binding,border,bottom,bottomright,bottomleft,break,caption-side,char,clear,clip,color,colors,cue,cursor,collapse,decoration,direction,display,elevation,empty-cells,family,filter,float,flow,focus,font,grid,height,image,increment,indent,input,inside,ime-mode,include-source,justify,last,layer,left,letter,line,list,margin,marker,marks,max,min,mode,modify,moz,offsett,opacity,orphans,outline,overflow,overflow-X,overflow-Y,overhang,position,position-x,position-y,padding,page,pause,pitch,play-during,position,quotes,radius,range,repeat,replace,reset,richness,right,ruby,set-link-source,select,shadow,size,spacing,speak,speak-header,speak-numeral,speak-punctuation,speech-rate,stress,stretch,style,table-layout,text,transform,text-autospace,text-kashida-space,top,topleft,topright,type,underline,unicode-bidi,use-link-source,user,variant,vertical,visibility,voice-family,volume,white-space,weight,widows,width,word,wrap,word-wrap,writing-mode,z-index,zoom}
}

% link URL, break long URL etc. -------------------------------------
\usepackage{url}


% include biblatex -----------------------------------------------------------------------
\usepackage[
	citestyle=my-authoryear-ibid,
	bibstyle=authoryear,
	backend=biber,
	babel=hyphen,
	urldate=comp,
	dateabbrev=false,
	maxcitenames=1
]{biblatex}
\usepackage[english=british]{csquotes}


% PDF Options -------------------------------------------------------------------------
\usepackage[
    bookmarks,
    bookmarksopen=true,
    colorlinks=true,
% use this colour definitions for links with colour
    linkcolor=red, % simple internal link
    anchorcolor=black,% anchor text
    citecolor=blue, % Anchor to bibliography entries in the text
    filecolor=magenta, % Anchors to local files
    menucolor=red, % Acrobat menu entries
    urlcolor=cyan, 
%
% use this colour definitions for pring: everything is black
%
    %linkcolor=black, % simple internal link
    %anchorcolor=black, % anchor text
    %citecolor=black, %Anchor to bibliography entries in the text
    %filecolor=black, % Anchors to local files
    %menucolor=black, % Acrobat menu entries
    %urlcolor=black, 
%
    plainpages=false, % for the correct creation of bookmarks
    pdfpagelabels, % for the correct creation of bookmarks
    hypertexnames=false, % for the correct creation of bookmarks
    linktocpage % link the page numbers instead of the whole toc entry
]{hyperref}

% Commands with print special characters cause errors if they are used with hyperref
\hypersetup{
    pdftitle={\myTitle \mySubtitle},
    pdfauthor={\myAuthor},
    pdfcreator={\creator},
    pdfsubject={\myTitle \mySubtitle},
    pdfkeywords={\keywords},
}

% For the glossary and abbreviations -------------------------
\usepackage[
	%xindy,
	nonumberlist,	% keine Seitenzahlen anzeigen
	acronym,			% ein Abkürzungsverzeichnis erstellen
	toc,					% Einträge im Inhaltsverzeichnis
	shortcuts			% acronym shortcuts
]{glossaries}

% Package for a clean inclusion of external PDF files -----------------
\usepackage[final]{pdfpages}

% Continuous numeration of footnotes -------------------------------
\usepackage{chngcntr}


% nicer tables --------------------------------------------------------------------
\usepackage{tabularx}
\usepackage{multirow}


% for long tables ---------------------------------------------------------------------
\usepackage{longtable}
\usepackage{ltxtable}
\usepackage{filecontents}
\usepackage{array}
\usepackage{ragged2e}
\usepackage{lscape}


% rotation of elements -------------------------------------------------------
\usepackage{rotating}


% Change the formatting of lists --------------------------------------------------
\usepackage{paralist}
\usepackage{enumitem}


% Necessary for the creation of own commands -------------------------------------
\usepackage{ifthen}
\usepackage{forloop}


% defines commands for \todo and \listoftodos --------------------------------
\usepackage[english]{todonotes}

% Enables strikethrough for text
\usepackage{ulem}


% used for the end of macros without parameter
\usepackage{xspace}

% used for bold labels in bibliography entries
\usepackage{xpatch}

% for multiple columns and rows
\usepackage{multicol}
\usepackage{multirow}

% better hyphenation etc.
\usepackage{microtype}

% This is used for the dummy text
\usepackage{blindtext}

%% Uncomment this, if you want to let chapters start at the current position and not on a new page. ------------------------------------------------------
%\usepackage{etoolbox}
%\makeatletter
%\patchcmd{\chapter}{\if@openright\cleardoublepage\else\clearpage\fi}{}{}{}
%\makeatother


% Erstellung eines Index und Abkürzungsverzeichnisses/Glossars aktivieren ------------------------------
\makeindex{}
%\makenomenclature{}


% Kopf- und Fußzeilen, Seitenränder etc. ---------------------------------------------------------------
% !TEX encoding = UTF-8 Unicode
% !TEX root =  Paper.tex

% Line spacing 1,5 lines ---------------------------------------------------------------------
\onehalfspacing{}


% Margins ------------------------------------------------------------------------------------
\setlength{\topskip}{\ht\strutbox} % behebt Warnung von geometry
\geometry{paper=a4paper,left=25mm,right=40mm,top=30mm}

\recalctypearea

% Font: Times New Roman
\renewcommand{\rmdefault}{ptm}
\renewcommand{\sfdefault}{ptm}
\renewcommand{\ttdefault}{cmtt}


% Head- and footline --------------------------------------------------------------------------
\pagestyle{scrheadings}

% pring head- and footline at the beginning of a chapter
\renewcommand*{\chapterpagestyle}{scrheadings} 

% font style of the headline
\renewcommand{\headfont}{\normalfont}


% Headline for onesided documents:
\ihead{}
\chead{}
\rohead{\pagemark}
%\setlength{\headheight}{21mm} % Height of the headline

% expand the headline further the text
%\setheadwidth[0pt]{textwithmarginpar} 
%\setheadsepline[text]{0.4pt} % seperation line beneath the headline

%% headline for twosided documents:
%\ihead{\textit{\leftmark}}
%\chead{}
%\ohead{}
%\lehead{\includegraphics[scale=0.25]{\logo}}
%\rohead{\includegraphics[scale=0.25]{\logo}}
%\setheadsepline[text]{0.4pt} % seperation line beneath the headline


% Footline
%\cfoot{-- DRAFT -- \\[4ex]} % uncomment during draft status
\cfoot{} % comment during draft status
\ofoot{}

% Remove the paragraph indentation
\setlength{\parindent}{0pt}


% 	miscellaneous settings ---------------------------------------------------
% a little bit more spacing after a fullstop
\frenchspacing{}

% avoide orphans and widow lines
\clubpenalty = 10000
\widowpenalty = 10000 
\displaywidowpenalty = 10000

% format source code
\lstset{numbers=left, numberstyle=\tiny, numbersep=5pt, breaklines=true}
\lstset{emph={square}, emphstyle=\color{red}, emph={[2]root,base}, emphstyle={[2]\color{blue}}}

% footnotes are counted continuously
%\counterwithout{footnote}{chapter}

% more space in tables
\renewcommand{\arraystretch}{\TableCellPadding}

% Changes to the Citation style
\renewcommand*{\postnotedelim}{\addcolon\space}

\DeclareFieldFormat{postnote}{#1}
\DeclareFieldFormat{multipostnote}{#1}

% Slashes instead of 'ands'
\let\oldmultinamedelim\multinamedelim
\let\oldfinalnamedelim\finalnamedelim
\renewcommand*{\multinamedelim}{\slash}
\renewcommand*{\finalnamedelim}{\slash}
\AtBeginBibliography{%
  \renewcommand*{\multinamedelim}{\oldmultinamedelim}%
  \renewcommand*{\finalnamedelim}{\oldfinalnamedelim}%
}

% Layout for bibliography
\setlength\bibitemsep{3\itemsep} % Set the space between bibliography entries
\setlength{\bibhang}{1cm}

% Change 'and' to ';'
\AtBeginBibliography{%
  \renewcommand*{\multinamedelim}{;\space}%
  \renewcommand*{\finalnamedelim}{;\space}%
}

% Bibliography Layout
\renewcommand*{\labelnamepunct}{\addcolon\space}
\renewcommand*{\bibpagespunct}{\addcolon\space}

\DeclareLanguageMapping{british}{custom-english-ordinal-sscript}

\DeclareFieldFormat[book,article,inbook,booklet,manual]{title}{#1}
\DeclareFieldFormat[inbook]{booktitle}{#1}
\DeclareFieldFormat{journaltitle}{\mkbibemph{#1}} % italic journal title with comma
\DeclareFieldFormat[article,inbook]{pages}{#1}

\DeclareFieldFormat{edition}%
                   {\ifinteger{#1}%
                    {\mkbibordedition{#1}\addnbthinspace{}edition}%
                    {#1\isdot}}
                    
\DeclareFieldFormat{url}{[\url{#1}}
\DeclareFieldFormat{urldate}{\addcomma\space\bibstring{urlseen}\space#1]}
\DeclareFieldFormat{editor}{#1}%\space\bibstring{byeditor}}

\renewbibmacro*{journal+issuetitle}{%
  \usebibmacro{journal}%
  \setunit*{\addspace}%
  \iffieldundef{series}
    {}
    {\newunit
     \printfield{series}%
     \setunit{\addspace}}%
    \usebibmacro{volume+year+issue}%
    \setunit{\addcolon\space}%
    \usebibmacro{issue}%
    \newunit}


\newbibmacro*{volume+year+issue}{%
	\printfield{volume}%
	\iffieldundef{number}%
		{}%
		{\setunit*{\adddot}%
		\printfield{number}%
		}%
	\printtext[parens]{%
  	\iffieldundef{year}%
  		{}%
  		{\printfield{year}}%
  	}%
  	\printtext{%
    	\iffieldundef{issue}%
      	{\usebibmacro{date}}%
      	{\printfield{issue}%
       	}}%
  \newunit}

%% Inbooks
\DeclareBibliographyDriver{inbook}{%
  \usebibmacro{bibindex}%
  \usebibmacro{begentry}%
  \usebibmacro{author/translator+others}%
  \setunit{\labelnamepunct}\newblock
  \usebibmacro{title}%
  \newunit
  \printlist{language}%
  \newunit\newblock
  \usebibmacro{byauthor}%
  \newunit\newblock
  \usebibmacro{in:}%
  \usebibmacro{bybookauthor}%
  \setunit{;\space}\newblock
  \usebibmacro{byeditor+others}%  
  \newunit\newblock
  \usebibmacro{maintitle+booktitle}%
  \newunit\newblock
  \printfield{edition}%
  \newunit
  \iffieldundef{maintitle}
    {\printfield{volume}%
     \printfield{part}}
    {}%
  \newunit
  \printfield{volumes}%
  \newunit\newblock
  \usebibmacro{series+number}%
  \newunit\newblock
  \printfield{note}%
  \newunit\newblock
  \usebibmacro{publisher+location+date}%
  \newunit\newblock
  \usebibmacro{chapter+pages}%
  \newunit\newblock
  \iftoggle{bbx:isbn}
    {\printfield{isbn}}
    {}%
  \newunit\newblock
  \usebibmacro{doi+eprint+url}%
  \newunit\newblock
  \usebibmacro{addendum+pubstate}%
  \setunit{\bibpagerefpunct}\newblock
  \usebibmacro{pageref}%
  \newunit\newblock
  \iftoggle{bbx:related}
    {\usebibmacro{related:init}%
     \usebibmacro{related}}
    {}%
  \usebibmacro{finentry}}

% Bold bibliography entries
%\xpretobibmacro{author}{\mkbibbold\bgroup}{}{}
%\xapptobibmacro{author}{\egroup}{}{}
%\xpretobibmacro{bbx:editor}{\mkbibbold\bgroup}{}{}
%\xapptobibmacro{bbx:editor}{\egroup}{}{}
%\renewcommand*{\labelnamepunct}{\mkbibbold{\addcolon\space}}


% eigene Definitionen für Silbentrennung ---------------------------------------------------------------
% !TEX encoding = UTF-8 Unicode
% !TEX root =  Paper.tex

% Trennvorschläge im Text werden mit \" angegeben
% untrennbare Wörter und Ausnahmen von der normalen Trennung können in dieser
% Datei mittels \hyphenation definiert werden




% eigene LaTeX-Befehle ---------------------------------------------------------------------------------
% !TEX encoding = UTF-8 Unicode
% !TEX root =  Paper.tex
% Custom commands


% simple changing of the font, e.g.: \changefont{cmss}{sbc}{n} ---------------------------------------
\newcommand{\changefont}[3]{\fontfamily{#1} \fontseries{#2} \fontshape{#3} \selectfont}

\newcommand{\bs}{$\backslash$}


% List elements with bold heading ------------------------------------------------------
\newcommand{\itemd}[2]{\item{\textbf{#1}}\\{#2}}


% pring authors
\newcommand{\AuthorName}[1]{\textsc{#1}}
\newcommand{\Author}[1]{\AuthorName{\citeauthor{#1}}}


% commands for highlighting words -----------------------------------------------
\newcommand{\NewTerm}[1]{\textbf{#1}}


% Amount with currency -----------------------------------------------------------------------------------
\newcommand{\Amount}[2][general]{#2\,\ifthenelse{\equal{#1}{dollar}}{\$}{}\ifthenelse{\equal{#1}{euro}}{€}{}\ifthenelse{\equal{#1}{yen}}{¥}{}\ifthenelse{\equal{#1}{cent}}{¢}{}\ifthenelse{\equal{#1}{pound}}{£}{}\ifthenelse{\equal{#1}{peso}}{₱}{}\ifthenelse{\equal{#1}{baht}}{฿}{}\ifthenelse{\equal{#1}{franc}}{₣}{}\ifthenelse{\equal{#1}{lira}}{₤}{}\ifthenelse{\equal{#1}{drachma}}{₯}{}\ifthenelse{\equal{#1}{pfennig}}{₰}{}\ifthenelse{\equal{#1}{general}}{¤}{}}


% Miscellaneous ---------------------------------------------------------------------------------------------
\newcommand{\Input}[1]{\texttt{#1}}
\newcommand{\Code}[1]{\texttt{#1}}
\newcommand{\File}[1]{\texttt{#1}}

\newcommand{\Filetype}[1]{\textsf{#1}}
\newcommand{\XMLElement}[1]{\textsf{#1}}
\newcommand{\Webservice}[1]{\textsf{#1}}


% Column defintion right-aligned with defined width -------------------------------------------------
\newcolumntype{w}[1]{>{\raggedleft\hspace{0pt}}p{#1}}


% Left-aligned table columns with tabularx -------------------------------------------------------------
\newcolumntype{y}[1]{>{\RaggedRight\arraybackslash\hsize=#1\hsize}X}



% Das eigentliche Dokument -----------------------------------------------------------------------------
%   Der eigentliche Inhalt des Dokuments beginnt hier. Die einzelnen Seiten
%   und Kapitel werden in eigene Dateien ausgelagert und hier nur inkludiert.
% ------------------------------------------------------------------------------------------------------
\begin{document}


% auch subsubsections nummerieren ----------------------------------------------------------------------
\setcounter{secnumdepth}{3}
% Nummerierungsebenen im Inhaltsverzeichnis
\setcounter{tocdepth}{2}


% Deckblatt und Abstract ohne Seitenzahl ---------------------------------------------------------------
\ofoot{}
% !TEX encoding = UTF-8 Unicode
% !TEX root =  Paper.tex

\thispagestyle{plain}
\begin{titlepage}

\changefont{cmss}{bx}{n}

\begin{center}

\LARGE{\textbf{\myTitleTitlepage}}\\
\Large{\textbf{\mySubtitleTitlepage}}\\[4ex]

\Large{\textbf{\kind}}\\[1.5ex]

\changefont{cmss}{m}{n}
\large{\area{} \\ \study}\\[1ex]
\large{at the \university{} \place}\\[6ex]

by

\large{\textbf{\myAuthor}} \\[2ex]

October \myYear

\vspace{\fill}

% 2 Logos -------------------
%\begin{tabular}{p{50mm}cp{50mm}c}
%	\includegraphics[scale=0.16]{Logo.pdf} & \includegraphics[trim=0mm 100mm 0mm 100mm, scale=0.25]{Logo.pdf}
%\end{tabular}

% 1 Logo ----------------------
\includegraphics[width=0.5\textwidth]{Logo.pdf}

\vspace{\fill}

\normalsize{}
\renewcommand{\arraystretch}{1}
\begin{tabular}{ll}
Time needed for completion:  & \quad 12 Weeks\\
Student ID, Course: & \quad \studentID, \course\\
Company: & \quad \company, company location\\
Supervisor:  & \quad \firstreviewer\\
Second reviewer: & \quad \secondreviewer\\
Submitted in: & \quad \place, 4th October 2013
\end{tabular}
\renewcommand{\arraystretch}{\TableCellPadding}

\end{center}

\end{titlepage}


% Sperrvermerk ----------------------------------------------------------------------------
% !TEX encoding = UTF-8 Unicode
% !TEX root =  Paper.tex

\chapter*{Sperrvermerk}

\thispagestyle{empty}


Die vorliegende {\art} mit dem Titel \textit{\titel} beinhaltet interne und vertrauliche Informationen der \firma{}, sie ist daher mit einem Sperrvermerk versehen und wird ausschließlich zu Pr\"ufungszwecken im Studiengang {\studienbereich} der \hochschule{} \ort{} vorgelegt.

Jede Einsichtnahme und Ver\"offentlichung -- auch von Teilen der Arbeit oder in digitaler Form –- bedarf der vorherigen schriftlichen Zustimmung durch die {\firma}.

% Selbständigkeitserklärung ----------------------------------------------------------------------------
% !TEX encoding = UTF-8 Unicode
% !TEX root =  Paper.tex

\chapter*{Eidesstattliche Versicherung}

\thispagestyle{empty}

Ich, \autor, Matrikel-Nr.\ \matrikelnr, versichere hiermit, dass ich meine \art\xspace mit dem Thema
\begin{quote}
\textit{\titel} \textit{\untertitel}
\end{quote}
selbständig verfasst und keine anderen als die angegebenen Quellen und Hilfsmittel benutzt habe, wobei ich alle wörtlichen und sinngemäßen Zitate als solche gekennzeichnet habe. Die Arbeit wurde bisher keiner anderen Prüfungsbehörde vorgelegt und auch nicht veröffentlicht. \\[6ex]

\ort, den \today\\[3ex]


\rule[-0.2cm]{5cm}{0.5pt}

\textsc{\autor} 


%  Abstract ------------------------------------------------------------------------------------------
% !TEX encoding = UTF-8 Unicode
% !TEX root =  ../Paper.tex

\chapter*{Abstract}
\label{cha:Abstract}

\thispagestyle{empty}


Abstract

% Liste mit Todos -------------------------------------------------------------------------------------
\phantomsection
%\markboth{Zu erledigende Punkte}{Liste der noch zu erledigenden Punkte}
\listoftodos
\newpage



% Seitennummerierung -----------------------------------------------------------------------------------
%   Vor dem Hauptteil werden die Seiten in großen römischen Ziffern 
%   nummeriert.
% ------------------------------------------------------------------------------------------------------
\ofoot{\pagemark \\[4ex]}
\pagenumbering{Roman}

% Inhaltsverzeichnis --------------------------------------------------------------------------------------
\phantomsection{} % Sorgt für korrekte Aufnahme des Inhaltsverzeichnisses in das Inhaltsverzeichnis
\addcontentsline{toc}{chapter}{Contents}
\tableofcontents{}


% Abkürzungsverzeichnis --------------------------------------------------------------------------------
% !TEX encoding = UTF-8 Unicode
% !TEX root =  ../Paper.tex

\addchap{Abbreviations}
\label{cha:Abbreviations}

\begin{acronym}[YTMMM]
\setlength{\itemsep}{-\parsep}

\acro{AGPL}{Affero GNU General Public License}
\acro{API}{Application Programming Interface}
\acro{WYSIWYG}{What You See Is What You Get}
\end{acronym}

% Glossar ---------------------------------------------------------------------------------------------

\makeglossaries
% !TEX encoding = UTF-8 Unicode
% !TEX root =  ../Paper.tex

%
% Glossareintraege --> referenz, name, beschreibung
% Aufruf mit \gls{...}
%
\newglossaryentry{Glossareintrag}{name={Glossareintrag},plural={Glossareinträge},description={Ein Glossar beschreibt verschiedenste Dinge in kurzen Worten}}

\printglossary[style=altlist, title=Glossar]
\label{cha:Glossary}


% Restliche Verzeichnisse ------------------------------------------------------------------------------
\listoffigures{} % Abbildungsverzeichnis
\listoftables{} % Tabellenverzeichnis
\renewcommand{\lstlistlistingname}{List of Listings}
\lstlistoflistings{} % Listings-Verzeichnis

% arabische Seitenzahlen im Hauptteil ------------------------------------------------------------------
\clearpage{}
\pagenumbering{arabic}


% die Inhaltskapitel werden in "Inhalt.tex" inkludiert -------------------------------------------------
% !TEX encoding = UTF-8 Unicode
% !TEX root =  Paper.tex

% Here you can include the different chapters. 
% They have to be in the corresponding .TEX file.
% If you want you can change the file names.


% Put all your chapter/section related text HERE (just copy paste)

\chapter{Introduction}
\label{cha:Introduction}

This is a simple introduction to the template. It shows several functions supported by the template like a glossary, acronyms, citing, referencing, tables and figures (see~\cref{sec:TablesFigures}), lists, listings, etc.

\section{Example of the Glossary}
\label{sec:ExGlossary}

This work is published under the \gls{CC-SA} license. So, feel free to use it for your thesis and adapt it to your needs. Because of the \ac{CC} license your allowed to do so. But first of all you should learn {\LaTeX}. Some people think its difficult to learn, because there is no \ac{WYSIWYG} editor.


\section{Citing}
\label{sec:Citing}

This template uses the Harvard citing style \enquote{author-year}. But it's to change, e.g. if you don't like in line citing with parentheses just start using footnotes. Here is an example; a wise quote by \textcite{Asimov1997}:

\begin{quote}
\enquote{Knowledge is indivisible. When people grow wise in one direction, they are sure to make it easier for themselves to grow wise in other directions as well. On the other hand, when they split up knowledge, concentrate on their own field, and scorn and ignore other fields, they grow less wise \textemdash\ even in their own field.}
\end{quote}

As you have guessed, this template doesn't use the default {\LaTeX} fonts. In order to do so, the package \emph{fontspec} is included. With \emph{fontspec} the user is able to load OpenType fonts in a {\LaTeX} document \autocite[see][3]{Robertson2013}. But in this case, \emph{Times New Roman} is chosen which is supported by {\LaTeX}.

\section{Tables and Figures}
\label{sec:TablesFigures}

This section has examples of tables (see~\cref{tab:SampleTable}) and figures (see~\cref{fig:PlaceHolder}).

\begin{figure}%
\caption{A nice and easy to use place holder.}%
\missingfigure{This is a place holder for a figure.}%
\caption*{Source: The reference is below the figure}
\label{fig:PlaceHolder}%
\end{figure}

\begin{center}
\begin{longtable}{c|c|c|c|c}
\caption{A sample table with performance data.}
\label{tab:SampleTable} \\

\textbf{\#}	&	\textbf{C} {[}sec{]}	&	\textbf{Java} {[}sec{]}	&	\textbf{Java - C} {[}sec{]}	&	\%	\\
\toprule[1.5pt]
\endfirsthead

\multicolumn{5}{c}{\tablename\ \thetable\ -- \textit{Continued from previous page}} \\
\textbf{\#}	&	\textbf{C} {[}sec{]}	&	\textbf{Java} {[}sec{]}	&	\textbf{Java - C} {[}sec{]}	&	\%	\\
\toprule[1.5pt]
\endhead
			
\multicolumn{5}{r}{\footnotesize \textit{Continued on next page}} \\
\endfoot
\endlastfoot

1	&	2.47	&	3.336932371	&	0.866932371	&	35.09847656	\\
2	&	2.49	&	3.28858181	&	0.79858181	&	32.07155863	\\
3	&	2.48	&	3.29198301	&	0.81198301	&	32.7412504	\\
4	&	2.49	&	3.286295345	&	0.796295345	&	31.97973273	\\
5	&	2.49	&	3.288482704	&	0.798482704	&	32.06757847	\\
6	&	2.48	&	3.290081162	&	0.810081162	&	32.66456298	\\
7	&	2.48	&	3.285896621	&	0.805896621	&	32.49583149	\\
8	&	2.49	&	3.289358723	&	0.799358723	&	32.10275996	\\
9	&	2.48	&	3.288659403	&	0.808659403	&	32.60723399	\\
10	&	2.49	&	3.335017952	&	0.845017952	&	33.93646394	\\
11	&	2.48	&	3.287597187	&	0.807597187	&	32.5644027	\\
12	&	2.49	&	3.293863277	&	0.803863277	&	32.28366574	\\
13	&	2.48	&	3.291244927	&	0.811244927	&	32.71148899	\\
14	&	2.48	&	3.286065776	&	0.806065776	&	32.50265226	\\
15	&	2.49	&	3.288609327	&	0.798609327	&	32.07266373	\\
16	&	2.49	&	3.288866762	&	0.798866762	&	32.08300249	\\
17	&	2.48	&	3.292376565	&	0.812376565	&	32.75711956	\\
18	&	2.49	&	3.290268407	&	0.800268407	&	32.13929345	\\
19	&	2.48	&	3.289624469	&	0.809624469	&	32.64614794	\\
20	&	2.48	&	3.288649694	&	0.808649694	&	32.6068425	\\
21	&	2.49	&	3.293171429	&	0.803171429	&	32.25588068	\\
22	&	2.48	&	3.291608451	&	0.811608451	&	32.72614722	\\
23	&	2.48	&	3.282929763	&	0.802929763	&	32.37620012	\\
24	&	2.49	&	3.293829963	&	0.803829963	&	32.28232783	\\
25	&	2.48	&	3.285785784	&	0.805785784	&	32.49136226	\\
\multicolumn{5}{c}{\ldots} \\
\bottomrule[2pt]
AVG:	&	2.4838	&	3.290900608	&	0.807100608	&	32.49540061	\\
STD	&	0.005674864	&	0.009707116	&	0.011960919	&	0.530352252	\\
\caption*{Source: The source of this table.}
\end{longtable}
\end{center}

\section{Lists}
\label{sec:Lists}

Lists look like this:

\blinditemize[5]


\chapter{How To}
\label{cha:HowTo}

This chapter is a guide on how to use this template. For example, the referencing and citing is explained. Moreover, the usage of figures and tables is elaborated.

\section{Build the document}
\label{sec:BuildDocument}

In order to build your document, you just have to double-click the file \File{makefile.bat}.

\section{Use of the Glossary}
\label{sec:UseGlossary}

In order to use the glossary, an entry has to be put in the file \File{Content/Glossary.tex}. \Cref{lst:glossary} shows a sample entry and the command for using it in the text. In the Glossary, a dot is automatically added at the end of each description.\footnote{This is an example of a footnote. Each page starts by 1 (can be changed for continuous numbering.}

\begin{lstlisting}[language={[LaTeX]TeX},label={lst:glossary},caption={The use of the Glossary.}]
\newglossaryentry{CC-SA}{
	name={Creative Commons Attribution-ShareAlike 3.0 Unported},
	description={This \acf*{CC} license allows the user to \emph{share} (copy and redistribute) a work and to \emph{adapt} (remix, transform and build upon) the work for \emph{any} purpose \autocite{creativecommons2014}}
}

@\ldots@ \gls{CC-SA} is used @\ldots@
\end{lstlisting}

\section{Use of Abbreviations/Acronyms}
\label{sec:UseAcronyms}

Acronyms/Abbreviations are also added to the file \File{Content/Glossary.tex}. \Cref{lst:acronyms} shows a sample acronym entry and the usage of it. There are many commands to use acronyms:

\begin{itemize}
	\item \lstinline[language={[LaTeX]TeX}]!\ac{<label>}! is the default command. First, the whole word is written. Afterwards only the short form is used
	\item \lstinline[language={[LaTeX]TeX}]!\acl{<label>}! always uses the long form
	\item \lstinline[language={[LaTeX]TeX}]!\acs{<label>}! is the short form
	\item \lstinline[language={[LaTeX]TeX}]!\acf{<label>}! writes the full acronym: Longform (shortform)
	\item Each command can be used with a \lstinline!*!, so that the produced text is not referenced\footnote{Another random footnote without important content.}
\end{itemize}

\begin{lstlisting}[language={[LaTeX]TeX},label={lst:acronyms},caption={The use of Abbreviations/Acronyms.}]
\newacronym{CC}{CC}{Creative Commons}
\newacronym{WYSIWYG}{WYSIWYG}{What You See Is What You Get}

\newacronym{<label>}{Shortform}{Longform}
\end{lstlisting}

\section{Referencing Commands}
\label{sec:ReferencingCommands}

Each figure, formulae or table can be accessed by a label for referencing. Therefore, two important commands are available:

\begin{enumerate}
	\item \lstinline[language={[LaTeX]TeX}]!\Cref{<label>}!
	\item \lstinline[language={[LaTeX]TeX}]!\cref{<label>}!
\end{enumerate}

The first capitalizes the first letter of \enquote{Figure}, \enquote{Tables}, etc. \Cref{lst:referencing} shows these two sentences in {\LaTeX} code where \texttt{lst:referencing} is the label of the listing.

\begin{lstlisting}[language={[LaTeX]TeX},label={lst:referencing},caption={How to reference.}]
The first capitalizes the first letter of \enquote{Figure}, \enquote{Tables}, etc. \Cref{lst:referencing} shows these two sentences in {\LaTeX} code where \texttt{lst:referencing} is the label of the listing.
\end{lstlisting}

\section{Quoting}
\label{sec:Quoting}

Quoting using this {\LaTeX} template is highly automated. There are several commands introduced in this section.

\begin{itemize}
	\item \lstinline[language={[LaTeX]TeX}]!\enquote{Text}!: Puts the text in quotation marks
	\item \lstinline[language={[LaTeX]TeX}]!\enquoteit{Text}!: Puts the text in quotation marks and italicizes the text (given by the Guidelines)
	\item \lstinline[language={[LaTeX]TeX}]!\foreignquote{<lang>}{Text}!: Should be used for quotes in foreign language, e.\,g.\,German.
	\item \lstinline[language={[LaTeX]TeX}]!\foreignquoteit{<lang>}{Text}!: See \lstinline[language={[LaTeX]TeX}]!\enquoteit{Text}!
	\item \lstinline[language={[LaTeX]TeX}]!\textquote[<cite>][<punct>]{<text>}<tpunct>!: This command can be used for direct quotations. It puts the \texttt{<text>} in quotation marks followed by a reference to the author. For \texttt{<cite>} you use the regular \lstinline[language={[LaTeX]TeX}]!\cite[<e.g. see>][<page>]{<citekey>}! command (in curly braces, not autocite command needed, because parenthesis are added automatically). \texttt{<punct>} is the punctuation after the text. \texttt{<tpunct>} is a punctuation after the whole command
\end{itemize}

The commands \lstinline[language={[LaTeX]TeX}]!\autocite[<e.g. see>][<page>]{<citekey>}! and \lstinline[language={[LaTeX]TeX}]!\textcite[<e.g. see>][<page>]{<citekey>}! are used for indirect quotes as usual.

\subsection{Examples}
\label{sec:QuotingExamples}

\begin{itemize}
	\item \lstinline[language={[LaTeX]TeX}]!\enquote{Text}!: \enquote{This text is in quotation marks}
	\item \lstinline[language={[LaTeX]TeX}]!\enquoteit{Text}!: \enquoteit{This text is in quotation marks and italicized}
	\item \lstinline[language={[LaTeX]TeX}]!\foreignquote{<lang>}{Text}!: \foreignquote{ngerman}{Dieser Text ist in Deutsch}
	\item \lstinline[language={[LaTeX]TeX}]!\foreignquoteit{<lang>}{Text}!: \foreignquoteit{ngerman}{Dieser Text ist in Deutsch und kursiv}
	\item \lstinline[language={[LaTeX]TeX}]!\textquote[<cite>][<punct>]{<text>}<tpunct>!: \textquote[{\cite[see][3]{Robertson2013}}]{Without fontspec, it is necessary to write cumbersome font definition files for {\LaTeX}, since \LaTeX{'s} font selection scheme (known as the \enquote{\textsc{nfss}} has a lot going on behind the scenes to allow easy commands like \texttt{\bs emph} or \texttt{\bs bfseries}}
	\item Multiple citations from the same author in a row are shorted: \cite[see][3]{Robertson2013}
	\item But only if they are on the same page: \cite[see][3]{Robertson2013}
\end{itemize}


\chapter{Some Blindtext}
\label{cha:SomeBlindtext}

\blindtext

\section{Some Paragraphs}
\label{sec:SomeParagraphs}

\Blindtext


\clearpage{}
\pagenumbering{Roman}
\setcounter{page}{7} %%% Dieser Pagecounter muss entsprechend der verbrauchten Seiten im Inhaltsverzeichnis angepasst werden. Endet das IHV bei Seite III, so muss hier 4 eingetragen werden


% Literaturverzeichnis ---------------------------------------------------------------------------------
%   Das Literaturverzeichnis wird aus der BibTeX-Datenbank "Bibliographie.bib"
%   erstellt.
% ------------------------------------------------------------------------------------------------------
\bibliography{Bibliography} % Aufruf: bibtex Masterarbeit
\bibliographystyle{agsm} % Harvard-Stil des Literaturverzeichnisses


% Index ------------------------------------------------------------------------------------------------
%   Zum Erstellen eines Index, die folgende Zeile auskommentieren.
% ------------------------------------------------------------------------------------------------------
%\printindex


% Anhang -----------------------------------------------------------------------------------------------
%   Die Inhalte des Anhangs werden analog zu den Kapiteln inkludiert.
%   Dies geschieht in der Datei "Anhang.tex".
% ------------------------------------------------------------------------------------------------------
\begin{appendix}
    \clearpage{}
    \pagenumbering{roman}
    \chapter{Appendix}
    \label{sec:Appendix}
    % Rand der Aufzählungen in Tabellen anpassen
    \setdefaultleftmargin{1em}{}{}{}{}{}
    % !TEX encoding = UTF-8 Unicode
% !TEX root =  Paper.tex

\section{CD Attached}
\label{sec:CDAttached}

\subsection{Contents of the CD}
\label{subsec:ContentsCD}

\begin{enumerate}
	\item The thesis as PDF
	\item All online resources as PDF
	\item Source code
	\item \ldots
\end{enumerate}
\end{appendix}

\end{document}
